\labsection{2}{ERP Benefits, Challenges, and Modular Integration}{sec:lab02}

\begin{labproject}{Teach one ERP benefit and one ERP challenge using the system you just built}
\textbf{Coverage.} Tutorial 2 + Chapter 2.\
\textbf{Core system.} Odoo 19 Community app structure.\
\textbf{Goal.} Move from generic statements such as ``ERP improves efficiency'' to evidence-based explanations of how integration creates benefits and why implementation is difficult.

\begin{labmode}
Tutorial 2 uses Team Benefit and Team Challenge. Odoo is used here as a concrete reference architecture: each team must point to at least one app-to-app connection or setup choice that demonstrates the selected benefit/challenge.
\end{labmode}

\medskip\noindent\textbf{Part A --- Team Benefit.}

Choose one ERP benefit such as shared data, faster process flow, management visibility, reduced duplicate entry, better inventory accuracy, or integrated reporting. Explain:
\begin{enumerate}
  \item the old fragmented problem;
  \item the ERP mechanism that changes it;
  \item the Odoo record or connection that demonstrates the mechanism;
  \item one metric that could show the benefit.
\end{enumerate}

Example structure: ``A Sales Order creates a linked Delivery and can create an Invoice. This reduces repeated entry of product, quantity, customer, and price. Measure the percentage of invoices requiring manual correction.''

\medskip\noindent\textbf{Part B --- Team Challenge.}

Choose one challenge such as data migration, training, change resistance, scope creep, poor process design, weak sponsorship, role/permission design, or excessive customisation. Explain:
\begin{enumerate}
  \item why it is difficult;
  \item which users/processes are affected;
  \item what could happen if it is ignored;
  \item how the implementation team should reduce the risk.
\end{enumerate}

\medskip\noindent\textbf{Part C --- Inspect modularity in Odoo.}

Open the app grid and identify which app owns each of these records: customer, opportunity, quotation, purchase order, receipt, BOM, manufacturing order, delivery, invoice, employee. Then draw arrows showing which records create or feed the next record.

\begin{keyidea}
The point is to see modularity without silos. The records live in different apps because the workflows differ, but they remain connected through shared master data and document relationships.
\end{keyidea}

\medskip\noindent\textbf{Part D --- Teach the other subgroup.}

Follow Tutorial 2: explain your selected topic, listen to the other subgroup, take notes, and write the required paragraph demonstrating what you learned.
\end{labproject}

\begin{labdeliverable}
Submit the paragraph required by Tutorial 2. Keep your research source details separately as required by the tutorial. In class, be ready to demonstrate one Odoo screen or linked-record example that makes your benefit/challenge concrete.
\end{labdeliverable}
